\chapter{Introduction}

Medical licensing examinations such as the \gls{usmle} represent one of the most demanding knowledge challenges in professional education. Candidates must synthesize clinical reasoning across dozens of specialties under timed, high-stakes conditions. Traditional preparation relies on self-directed study, peer quizzing, and expensive tutoring. These resources are not equally accessible to all candidates. 

This project presents \gls{emma}, a conversational medical study agent built to close that gap. In explanation mode, students pose clinical questions in natural language and receive responses grounded in passages from eighteen standard medical textbooks. In quiz mode, \gls{emma} presents authentic \gls{usmle}-style questions, evaluates answers, and tracks per-specialty performance over time. A collaborative filtering recommender then steers students toward their weakest areas.
\gls{emma} is built on a five-component \gls{ml} pipeline: a SpaCy \gls{ner} model for clinical entity extraction, a \gls{faiss} vectorstore over 36,723 textbook chunks, a \gls{tfidf} + \gls{lsvc} specialty classifier, a BERTopic clustering model for fine-grained topic routing, and a Surprise-based \gls{cf} recommender. All components feed a locally-run Qwen3-4B-Thinking-2507 \gls{llm}. The central empirical question is whether textbook-grounded \gls{rag} improves a small local \gls{llm}'s accuracy on MedQA compared to a no-retrieval baseline. A six-run ablation grid shows the answer is condition-dependent: with the right embedding and \gls{llm} pairing, \gls{rag} delivers an +11 \gls{pp} accuracy gain; without those conditions, it can hurt. This nuanced finding is the primary contribution of this report.